% Tableau des labs du cours SIF (SEC3266) dans NDAP SUKULU (propre au cours).
Le cours compte dix laboratoires, disponibles sous le code \codeUE. Le
dernier reconstitue l'attaque de l'Opération GAROUA à partir de ses
journaux.

\begin{tableaucours}{L{2.4cm}L{6.4cm}L{1.6cm}L{1.8cm}}{%
    \tableauHeader{Lab} &
    \tableauHeader{Sujet} &
    \tableauHeader{Chapitre} &
    \tableauHeader{Opération}
}
SIF-LAB-01 & Empreintes et HMAC~: retrouver les fichiers modifiés &
\ref{chap:crypto-sym} & GAROUA \\
SIF-LAB-02 & Certificats~: lire une chaîne, repérer les défauts &
\ref{chap:crypto-asym} & LIZBIZ \\
SIF-LAB-03 & Analyse de risques EBIOS simplifiée & \ref{chap:risques} &
LIZBIZ \\
SIF-LAB-04 & Classer les actifs et coter les besoins & \ref{chap:actifs} &
LIZBIZ \\
SIF-LAB-05 & Capture réseau~: clair, balayage, ARP & \ref{chap:reseaux1} &
GAROUA \\
SIF-LAB-06 & Règles de pare-feu & \ref{chap:reseaux2} & GAROUA \\
SIF-LAB-07 & Revue des comptes & \ref{chap:iam} & GAROUA \\
SIF-LAB-08 & Durcissement Linux & \ref{chap:systemes} & LIZBIZ \\
SIF-LAB-09 & Revue de code & \ref{chap:applications} & MEFOU \\
SIF-LAB-10 & Journaux~: reconstituer l'attaque & \ref{chap:operations} &
GAROUA \\
\end{tableaucours}
\vspace{-0.8em}
\begin{center}{\small\itshape Les laboratoires du cours dans NDAP SUKULU}\end{center}
